\documentclass{article}
\usepackage{graphicx} % Required for inserting images
\usepackage{amsmath}
\usepackage{amsfonts}
\usepackage{listings}
\usepackage{booktabs}
\usepackage{float}
\usepackage{hyperref} % should always go last
% \usepackage{enumitem}

\lstset{
  basicstyle=\ttfamily\footnotesize,
  frame=single,
  captionpos=t,
  breaklines=true,
}

\newcommand{\includesource}[1]{%
  \lstinputlisting[caption={\texttt{#1}}]{#1}%
}

\title{Performance Lab 6 - Group 8}
\author{Hans Bouwman, Joel de Haan}

\begin{document}

\maketitle


\tableofcontents

\newpage

\section{}

\section{Matrix-vector parallelisation}
To parallelise this routine a \texttt{\#pragma omp parallel schedule(static)}

With OpenMP the 3 scheduling policies are
\begin{itemize}
  \item static: assign work at compile time, threads know what they have to do. Advantage: work allocation doesn't happen at runtime. Disadvantage: if a thread stalls (goes slower than the rest) other threads in the pool can't help and just wait idle
  \item dynamic: assign each piece of work to a queue, the threads take the jobs from the queue. Advantage: minimises the chance of an idle thread. Disadvantage: while the queue can be made at compile time, popping tasks and assignation hapens at runtime
  \item guided: assign larger blocks of work at compile time. but still make threads take this pieces of work. In theory the best of both worlds but still has the overhead of a dynamical queue and stalling is still possible.
\end{itemize}

The choice depends strongly on the regularity of the workload. In this case, we are paralellising a row * column operation. Each iteration has exactly the same amount of operations (and most likely) a similar amount of memory acesses.
This is the most optimal kind of loop to allocate work statically.

After running a few runs in the cluster with the Kunpeng 920 nodes, the following results were obtained

Serial parallelisation:
    \begin{tabular}{|c|c}
          $N$ & time \\
         \hline
          1024 & 0.165 s \\
          2048 & 0.578 s \\
          4096 &  2.680 s \\
         8192 &  9.202 s  \\
         16384 & 42.931 s \\
         32768 & 138.412 s  \\
    \end{tabular}

Dynamic parallelisation:
    \begin{tabular}{|c|c}
          $N$ & time \\
         \hline
          1024 & 0.150 s \\
          2048 & 0.592 s \\
          4096 &  2.355 s \\
         8192 &  10.266 s  \\
         16384 & 39.884 s \\
         32768 &  155.946 s  \\
    \end{tabular}

The results were run on 100 cores. While it does show that the serial parallelisation is faster on the serial case for larger matrices/vectors
the results are, in general, very close to each other. This contradicts the fact that serial parallelisation has significantly less overhead.


\end{document}
